% Vista científica exacta materializada desde una rama condicional.
% Fuente: source/demostraciones_primarias/14_05b_extension_superviviente_caracter.tex; rama: HMTIncludeRouteLedgerBlock.
% No modifica el contenido matemático de la rama de origen.
\chapter{Extensión superviviente, registro dodecafásico y firma integral}
\label{chap:extension-ruta-caracter}

La trayectoria que se observa no constituye el objeto primario de esta
construcción. El objeto primario es una extensión compatible a todas las
profundidades del transporte HMT; una ruta finita es una de sus lecturas. Esta
distinción permite formular sin ambigüedad la cadena interna que conduce desde
una historia enriquecida hasta el carácter de acción. También determina qué
información puede olvidarse sin alterar la salida y cuál debe permanecer en el
registro.

El propósito de la sección es reunir en un solo argumento cuatro piezas que se
habían estudiado por separado: la extensión superviviente, la memoria
intrínseca, el registro orientado y la firma entera. El resultado terminal de
esta sección es el carácter formal
\[
 \Gamma_{108}^{\rm enr}
 \xrightarrow{\ L_{\rm tick}\ }
 \ZZ^5
 \xrightarrow{\ \chi_{\rm form}\ }
 \operatorname{Fun}(\mathcal P,\RR_{>0}),
\]
anterior tanto a la especialización de los coeficientes como a toda unidad
física. El primer mapa extrae la firma del registro leído; el segundo la
convierte en una función positiva sobre el espacio de parámetros
\(\mathcal P\). La especialización numérica sólo se efectúa después de
construir los coeficientes globales en la sección correspondiente. La
realización metrológica pertenece a una carta dimensional todavía posterior.
El mapa \(L_{\rm tick}\) es el lector integral de iteración y retorno del
registro cronológico enriquecido.

\section{Composición operatoria de la prolongación y cocadena torsional en la generación del registro de ruta}

La ley de masas comienza antes de la fórmula exponencial. APP aporta un único
soporte de rutas y sus dos hojas, suma y producto. El TRIT orienta cada
iteración elemental.
Si
\[
 r_t=1+((t-1)\bmod9),\qquad m_t:=M_{\rm ph}(r_t),
\]
el selector de fase \(M_{\rm ph}\) activa
\[
 \operatorname{op}_t=
 \begin{cases}
  \Sigma,&m_t=+1,\\
  \Pi,&m_t=-1,\\
  \varnothing\quad(\text{\rm retención o contador neutral}),&m_t=0.
 \end{cases}
\]
La lectura temporal puede interpretar el tercer caso como umbral o presente,
pero no lo convierte en una tercera hoja aritmética. No es APP quien alterna
sus hojas. TPK transporta
conjuntamente posición, orientación, hoja, fase, bloque, acarreo, prefijo,
cilindro, genealogía y memoria. El registro cronológico enriquecido anota
cada transición ---posición,
hoja, lectura, acarreo, fase, bloque y evento---. La firma local de frontera
se deriva después mediante \(\operatorname{Sig}_q(h,c,\lambda)\), y la firma
global de masa se extrae del registro completo; ninguna se reconstruye
retrospectivamente a partir de una masa.

El embudo finito que alimenta esta historia conserva cuatro tipos distintos:
\[
 9^2\!\cdot9^2\!\cdot4\!\cdot4=104\,976
 \longrightarrow 468\ \text{emisiones }U_6
 \longrightarrow 243\ \text{palabras }w_6
 \hookrightarrow \mathbb F_3^6,
 \qquad |\mathbb F_3^6|=729.
\]
La primera cifra cuenta condiciones iniciales del motor de dos cursores; 468
cuenta emisiones decimales y 243 palabras ternarias visibles. Las 324 semillas
de la proyección cronológica \(\Pi_{108}^{\rm cal}\) de un solo cursor
tampoco sustituyen las 104\,976
condiciones del estado completo. Cada flecha cambia, por tanto, de dominio y de
operación. La igualdad numérica entre el ambiente de \(729\) palabras y los
\(729\) sufijos de una etapa de filtración sólo se relaciona mediante la codificación
ternaria declarada; no identifica ambos objetos.

La elevación conserva asimismo su orden causal:
\[
 w_6\xrightarrow{\,L_0\,}w_{12}
 \xrightarrow{\,L_0\,}w_{18}
 \xrightarrow{\,L_1\,}w_{24}
 \xrightarrow{\,L_1\,}w_{30}
 \longrightarrow R_{36}\longrightarrow\mathcal G_9
 \longrightarrow
 \text{sección superviviente}.
\]
La frontera \(R_{36}\) abre la prolongación. En cada etapa de filtración,
\(\mathcal G_9\) agota
los 729 subcilindros y conserva el único compatible con todos los horizontes; el
episodio especular profundo \(3\to1\) es otra poda y no debe confundirse con
esa partición ordinaria \(729\to1\). El retorno nonádico conserva la fase,
pero cambia bloque, cilindro y memoria. Ni cifras objetivo ni masas intervienen
en esta selección. En la frontera dimensional, una partícula es una sección
persistente enriquecida de la extensión; la ruta es su historia observable.
Su nombre físico se lee después y no decide qué rama sobrevive.

Antes de abrir el registro cronológico dimensional, la misma genealogía conserva una
ventana publicada de veinte sextetos por canal y admite la primera lectura
arquimediana de doce tríadas decimales. Estos dos testigos pertenecen a la
prolongación del estado: no son los doce eventos del registro de masa, ni los
reemplazan. Los eventos se construyen sólo después, al leer la ruta sobre el
sistema dodecafásico orientado.

La genealogía dodecafásica convierte después los doce eventos orientados en un
registro y sólo entonces en la firma
\[
 \gamma_x\longmapsto\mathcal L(\gamma_x),\qquad
 \sigma_{\rm reg}:=L_{\rm tick}(\gamma_x)
 =(n_A,s_C,k_{\Dfour},b_{120},b_{\zeta}),
\]
\[
 \sigma_{\rm reg}\longmapsto
 \chi_{\rm form}(\sigma_{\rm reg};\,\cdot).
\]
Las coordenadas públicas \(b_{120}:=\nu_{120}^{\rm ev}\) y
\(b_{\zeta}:=\nu_{270}\) son salidas del registro cronológico, no correcciones
metrológicas introducidas al final. Los coeficientes posteriores del
semigrupo de torres conservan otra procedencia:
\[
 N_{120}=b_{120}+\eta_{120},\qquad
 b_{\zeta}\zeta_{270}\not\equiv\eta_{270}s_{270},
\]
con \(\zeta_{270}=135\Dfour^2\) y
\(s_{270}=q_-^{270}-q_+^{270}\). El refinamiento de torres se añade después
del carácter central sin borrar la procedencia de ninguna coordenada.

Tres salidas laterales conservan también su tipo. El sector nulo abandona la
rama positiva mediante su lector y su carta nulos; no es el límite
\(m\to0\) de un carácter positivo. Una trayectoria virtual es una extensión
finita con obstrucción terminal y conserva registro cronológico y acarreo sin adquirir por
ello una masa imaginaria. Finalmente, la sombra Witt acompaña cada etapa de filtración y
preserva la genealogía excepcional, pero no selecciona la ruta ni los
coeficientes de masa.

\begin{figure}[p]
\centering
\begin{tikzpicture}[
  font=\small,
  node distance=3.5mm,
  box/.style={draw,rounded corners=1.5mm,align=center,text width=12.8cm,
    inner xsep=6pt,inner ysep=4pt},
  hmt/.style={box,draw=hmtblue,fill=hmtlightblue},
  md/.style={box,draw=hmtviolet,fill=hmtlightviolet},
  act/.style={box,draw=hmtgreen,fill=hmtlightgreen},
  result/.style={box,draw=hmtgold,fill=hmtlightgold},
  arr/.style={-{Stealth[length=2.2mm]},line width=.7pt}]

  \node[hmt] (app) {\textbf{APP}: toro aritmético nonádico, caminos y hojas
    aditiva y multiplicativa. La especificación completa contiene
    \(104\,976\) semillas de doble cursor.};
  \node[hmt,below=of app] (trit) {\textbf{TRIT}: orientación local de cada
    transición; el levantamiento conserva residuo, cociente y acarreo.};
  \node[hmt,below=of trit] (tpk) {\textbf{TPK}: transporte de fase con memoria;
    \(104\,976\to468\,U_6\to243\,w_6\subset\mathbb F_3^6\).};
  \node[hmt,below=of tpk] (surv) {Filtración coinductiva y monodromía
    \(\mathcal G_9\): trayectoria superviviente con genealogía conservada.};
  \node[md,below=of surv] (route) {La trayectoria de partícula produce,
    sin consultar la masa, un registro dodecafásico y la firma
    \(\sigma=(n_A,s_C,k_{\Delta_4},b_{120},b_\zeta)\in\mathbb Z^5\).};
  \node[md,below=of route] (chi) {El carácter central
    \(\chi_0(\sigma)\) y el refinamiento armónico
    \(\eta\in\mathbb Z^{\langle90,120\rangle}\) construyen
    \(\chi_{\rm full}(\sigma,\eta)\in\mathbb R_{>0}\).};
  \node[act,below=of chi] (action) {Rama dimensional independiente:
    \(\Sigma_H=\varphi\alpha_{\rm HMT}^{16}\to\hbar_{\rm int}\),
    \(\omega_0=2\pi/(108t_0)\),
    \(m_0=\hbar_{\rm int}\omega_0c_{\rm int}^{-2}\).};
  \node[result,below=of action] (mass) {Confluencia tipada:
    \(m_X^{\rm int}=m_0\,\chi_{\rm full}(\sigma_X,\eta_X)\).
    La escala común fija la normalización absoluta y se cancela en
    \(m_X/m_Y\).};
  \node[result,below=of mass] (compare) {La carta de unidades publica la magnitud.
    CODATA y PDG se abren únicamente en la columna final de contraste.};

  \foreach \a/\b in {app/trit,trit/tpk,tpk/surv,surv/route,route/chi,
    chi/action,action/mass,mass/compare}{\draw[arr] (\a)--(\b);}
\end{tikzpicture}

\caption{Derivación de una magnitud de masa a partir de una trayectoria HMT.
El carácter aporta una razón adimensional; la escala de acción y el reloj
aportan la sección dimensional. La referencia externa interviene sólo después
de aplicar la carta de unidades.}
\label{fig:derivacion-masa-rev9}
\end{figure}

\section{Proyección de una extensión superviviente sobre su ruta observable}

Sea \((\mathcal S_m,\rho_{m+1,m})_{m\geq0}\) el sistema de estados HMT
supervivientes a profundidad \(m\), con mapas de truncamiento
\(\rho_{m+1,m}:\mathcal S_{m+1}\to\mathcal S_m\). Definimos
\[
 \Omega_{\HMT}^{\rm surv}
 :=\varprojlim_m\mathcal S_m
 =\left\{(x_m)_m:\rho_{m+1,m}(x_{m+1})=x_m\right\}.
\]
Un elemento \(x\in\Omega_{\HMT}^{\rm surv}\) conserva simultáneamente la
compatibilidad de todos sus truncamientos. Una carta de lectura de profundidad
108 produce
\[
 \operatorname{sh}_{108}(x)=\gamma_x\in\Gamma_{108}^{\rm enr}.
\]
La notación \(\operatorname{sh}\) recuerda que \(\gamma_x\) es una sombra
observable: registra una presentación secuencial del estado, pero el estado no
se define escogiendo retrospectivamente esa ruta.

La presentación enriquecida conserva, además de la palabra visible, fase,
hoja, orientación, borde, memoria antihoja y los predecesores necesarios para
evaluar los cociclos. Dos palabras con los mismos cinco enteros finales pueden
ser presentaciones distintas porque esos enteros se calculan después de
aplicar signos y correlaciones.

\begin{definition}[Memoria intrínseca y registro leído]
Una extensión dodecafásica enriquecida es una clase compatible de
truncamientos. Denotamos por \(\mathscr M(x)\) la memoria intrínseca que esos
truncamientos conservan. Fijada la carta
\(\gamma_x=\operatorname{sh}_{108}(x)\), su registro leído es
\[
 \mathcal L(\gamma_x)
 =\bigl(n_A,e_0,\ldots,e_{11},T_\zeta,C_\zeta\bigr),
\]
donde \(n_A\) es el conteo del soporte de acción declarado, los \(e_m\) son
los doce eventos orientados y \(T_\zeta,C_\zeta\) son la memoria torsional
total y su restricción a la corona. La compatibilidad exige que estos datos
sean preservados por todo truncamiento que contenga los pasos empleados en su
definición.
\end{definition}

La extensión ya contiene la memoria que la ruta presenta; la ruta no la crea.
Relativamente a la carta declarada, la composición
\[
x\longmapsto\gamma_x
\longmapsto\mathcal L(\gamma_x)
\longmapsto L_{\rm tick}(\gamma_x)
\]
es una aplicación. La invariancia frente a otra carta de presentación
requeriría, además, demostrar la compatibilidad de los correspondientes mapas
de transición. En este sentido preciso, el registro no se añade a una
partícula ya construida: es una lectura de la memoria de la extensión.

\section{Descomposición integral de \(12\) eventos en los sectores \(1+5+6\)}

Sea \(d_1,\ldots,d_{108}\in\{1,\ldots,9\}\) la palabra de una presentación
enriquecida. En la convención de paridad usada por el extractor, el soporte de
acción es
\[
 \mathcal D_{\rm act}=\{1,3,5,7\},
 \qquad
 n_A=\#\{t:d_t\in\mathcal D_{\rm act}\}.
\]
La elección de este soporte es parte declarada de la carta; no debe
confundirse con la clasificación TRIT de los pasos. Para los eventos pares se
usa
\[
 \mathcal D_{\rm evt}=\{2,4,6,8\},
 \qquad
 \Sigma_{12}=(+,+,+,-,-,-,+,+,+,-,-,-).
\]
El evento orientado del bloque \(m\in\{0,\ldots,11\}\) es
\[
 e_m=\Sigma_{12}(m)
 \#\{t\in\{9m+1,\ldots,9m+9\}:d_t\in\mathcal D_{\rm evt}\}.
\]

Emparejemos las dos semicoronas mediante
\[
 p_i=e_i+e_{i+6},
 \qquad
 a_i=e_i-e_{i+6},
 \qquad 0\leq i<6.
\]
Definimos
\[
 s_C=\sum_{i=0}^{5}p_i,
 \qquad
 M_i=p_i-p_5\quad(0\leq i<5),
\]
y escribimos \(M_5=(M_0,\ldots,M_4)\),
\(A_6=(a_0,\ldots,a_5)\).

El \cref{thm:memoria-1-5-6}, en la residencia primaria del sistema
dodecafásico de TPK, demuestra que
\[
 (e_0,\ldots,e_{11})\longmapsto(s_C,M_5,A_6)
\]
es inyectivo sobre su imagen y proporciona la inversa
\[
 p_5=\frac{s_C-\sum_{i=0}^{4}M_i}{6},\qquad
 p_i=p_5+M_i,
\]
\[
 e_i=\frac{p_i+a_i}{2},\qquad
 e_{i+6}=\frac{p_i-a_i}{2}.
\]
Aquí sólo se recuperan esas coordenadas para construir el carácter; la
prueba y la identidad \(1+5+6=12\) no se duplican.

El dato \(s_C\) es el que entra en el carácter angular. Las once coordenadas
\((M_5,A_6)\) conservan la colocación de ese total. Eliminarlas antes de
componer historias confunde eventos diferentes que pueden producir
correlaciones distintas.

\section{Cocadena torsional a nivel de paso}

El conteo de eventos no determina por sí solo la contribución torsional. Para
cada aparición de nueve, conservamos la cifra que la precede y definimos
\[
 j_2(d)=
 \begin{cases}
 0,&d=9,\\
 +1,&d\in\{2,4,6,8\},\\
 -1,&d\in\{1,3,5,7\}.
 \end{cases}
\]
Con índices cíclicos en la palabra, sea
\[
 \zeta_t=\mathbf 1_{\{d_t=9\}}j_2(d_{t-1}),
\]
\[
 T_\zeta=\sum_{t=1}^{108}\zeta_t,
 \qquad
 C_\zeta=\sum_{t\in\{27,54,81,108\}}\zeta_t.
\]
La coordenada fuerte que acompaña a \(\Dfour\) es
\[
 \boxed{k_{\Dfour}=T_\zeta-2C_\zeta.}
\]
La fórmula separa el arrastre acumulado de su corrección orientada en los
cuatro cierres de corona. El censo exhaustivo muestra que ni
\(T_\zeta\) ni \(C_\zeta\) por separado separan todas las fibras del censo;
el par ordenado sí lo hace en las dos convenciones tipadas que se contrastan.

Las dos coordenadas de torre restantes se calculan a partir del registro de
eventos. Sea
\[
 \tau_m=\operatorname{sgn}(e_m).
\]
Para las cuatro clases de paso cuatro,
\[
 I_a=\{a,a+4,a+8\}\pmod{12},
 \qquad 0\leq a<4,
\]
definimos el carácter ponderado por eventos
\[
 \nu_{120}^{\rm ev}
 =\sum_{a=0}^{3}\operatorname{sgn}
 \left(\sum_{m\in I_a}e_m\right),
\]
y la autocorrelación de paso tres
\[
 \nu_{270}=\sum_{m=0}^{11}\tau_m\tau_{(m+3)\bmod12}.
\]
La variante
\[
 \nu_{120}^{\rm sgn}
 =\sum_{a=0}^{3}\operatorname{sgn}
 \left(\sum_{m\in I_a}\tau_m\right)
\]
es otro observable. No se identifican silenciosamente: difieren en
25 de las 81 rutas del atlas.

\begin{definition}[Extractor enriquecido de firma]
Fijada la convención de eventos, se define
\[
 L_{\rm tick}(\gamma)=
 \bigl(n_A,s_C,k_{\Dfour},\nu_{120}^{\rm ev},\nu_{270}\bigr)
 \in\ZZ^5.
\]
El subíndice indica que la coordenada torsional se calcula antes de olvidar
los predecesores de los pasos nulos.
\end{definition}

Esta aplicación es determinista y no consulta masas, unidades ni etiquetas
de partículas. No se afirma que sea lineal respecto de una concatenación ya
proyectada: los signos y las autocorrelaciones obligan a componer primero los
registros enriquecidos.
